Deepfake detection research spans spatial, frequency, temporal, and multimodal cues, as well as dataset curation and evaluation protocols. Recent comprehensive surveys~\cite{verdoliva2020forensics_survey,pei2024deepfakesurvey,electronics2024survey} highlight two persistent gaps that motivate our work: (i)~detectors that perform well on curated benchmarks often degrade sharply under distribution shift, and (ii)~most evaluations assume server\hyp class hardware rather than mobile devices. These gaps underscore the need for system designs that treat cross\hyp dataset generalization and mobile constraints as first\hyp class objectives rather than afterthoughts.

\paragraph{Deepfake Detection as a Problem Space.}
Deepfake detection is fundamentally a binary classification task: given an image or video frame, determine whether it depicts a real or synthetically manipulated face. The challenge lies in the diversity of manipulation techniques (face swapping, reenactment, attribute editing, full synthesis) and the variability of real\hyp world conditions (compression artifacts, lighting, camera sensors, post\hyp processing). Early approaches relied on handcrafted forensic features such as head pose inconsistencies~\cite{yang2019exposing}, eye blinking patterns, or JPEG compression artifacts~\cite{cozzolino2019noiseprint}. Modern methods leverage deep learning to extract hierarchical representations directly from pixels or frequency spectra, achieving high accuracy on in\hyp distribution benchmarks~\cite{rossler2019faceforensicspp,dolhansky2020dfdc,li2020celebdf}.

However, cross\hyp dataset generalization remains elusive: models trained on FaceForensics++ or CelebDF can exhibit AUC drops of 20–30\% when evaluated on in\hyp the\hyp wild datasets like WildDeepfake~\cite{wilddeepfake2021} or Deepfake\hyp Eval\hyp 2024~\cite{deepeval2024}. This performance gap stems from dataset bias (e.g., training sets overrepresent specific generators or compression settings), domain shift (e.g., social media platforms apply proprietary compression pipelines), and the rapid evolution of synthesis methods. Addressing these gaps requires multi\hyp dataset training protocols that expose models to diverse artifacts, as well as rigorous cross\hyp dataset and out\hyp of\hyp distribution (OOD) evaluation protocols that measure robustness beyond in\hyp distribution accuracy.

\paragraph{Mobile Deployment Constraints.}
Most deepfake detection research assumes unlimited computational resources—models are evaluated on server GPUs with batch sizes and latencies that are impractical for on\hyp device inference. Yet many real\hyp world applications (e.g., live video verification, client\hyp side content moderation, edge forensics) require deployment on commodity smartphones or embedded devices with constrained memory, power budgets, and thermal envelopes. Mobile deployment introduces additional challenges: (i)~model size must fit within typical app bundles (50–100\,MB), (ii)~per\hyp frame latency must support real\hyp time or near\hyp real\hyp time inference (e.g., \textless200\,ms), and (iii)~energy consumption must preserve battery life across extended usage. These constraints motivate the use of lightweight architectures, model compression techniques, and adaptive inference strategies that balance accuracy and efficiency.

\paragraph{Core Detection Approaches.}
Modern deepfake detectors primarily use convolutional neural networks (CNNs) or vision transformers (ViTs) as feature extractors. CNN\hyp based approaches~\cite{rossler2019faceforensicspp,tan2019efficientnet,efficientnetv2} leverage hierarchical convolutional layers to capture spatial artifacts such as blending boundaries, facial inconsistencies, and generative noise patterns. The MobileNet~\cite{howard2017mobilenetv1,sandler2018mobilenetv2,howard2019mobilenetv3,mobilenetv4} and EfficientNet~\cite{tan2019efficientnet,efficientnetv2} families exemplify architectures optimized for mobile and resource\hyp constrained settings through depthwise separable convolutions and compound scaling.

Frequency\hyp domain methods~\cite{durall2020watch,qian2020thinkingfrequency,luo2024freqnet} exploit the observation that GAN\hyp generated images often exhibit distinctive patterns in their Fourier or DCT spectra, providing complementary cues that survive spatial\hyp domain compression. Recent work demonstrates that combining spatial and frequency features improves robustness to JPEG compression and other post\hyp processing operations~\cite{wang2022afd,zhang2024frequency_dcT}.

Vision transformers~\cite{dosovitskiy2021vit,liu2021swin} offer global receptive fields that can capture long\hyp range dependencies and subtle manipulation cues across entire frames. Hybrid CNN\hyp Transformer architectures~\cite{coccomini2022efficientnet_vit,khan2022hybrid_transformer_df,genconvit_github} aim to combine the inductive biases of CNNs (e.g., translation equivariance) with the modeling capacity of transformers, though at the cost of increased computational complexity.

Temporal modeling via 3D CNNs~\cite{tran2015c3d,carreira2017i3d}, LSTMs~\cite{zhao2022mobilenet_lstm,xception_convlstm2022}, or video transformers~\cite{bertasius2021timesformer} extends frame\hyp level detection to video sequences, enabling the use of inter\hyp frame consistency and motion artifacts as additional forensic signals. However, temporal approaches typically require multi\hyp frame context windows and are challenging to deploy on mobile devices due to memory and latency constraints.

\paragraph{Model Compression and Efficiency.}
To bridge the gap between research models and deployable systems, practitioners rely on three primary compression techniques:
\begin{itemize}[topsep=4pt,itemsep=2pt]
  \item \textbf{Pruning}~\cite{han2015learningboth,han2016deepcompression,blalock2020pruningsurvey}: Removing redundant weights or channels to reduce model size and FLOPs. Structured pruning (removing entire filters) offers better hardware efficiency than unstructured (weight\hyp level) pruning.
  \item \textbf{Quantization}~\cite{jacob2017integeronly,esser2020lsq,gholami2021quant_survey}: Representing weights and activations with lower\hyp precision data types (e.g., INT8 instead of FP32). Post\hyp training quantization (PTQ) is fast but limited to 8\hyp bit precision, while quantization\hyp aware training (QAT) enables more aggressive compression (e.g., INT4) at the cost of retraining.
  \item \textbf{Knowledge distillation}~\cite{hinton2015distilling,romero2015fitnets,gou2021kdsurvey}: Transferring knowledge from a large teacher model to a compact student model by matching soft probability distributions or intermediate feature representations.
\end{itemize}
Modern deep learning frameworks (PyTorch, TensorFlow, ONNX Runtime) provide end\hyp to\hyp end toolchains for converting research models into mobile\hyp optimized INT8 variants and exporting them to Android (via ONNX, TFLite) or iOS (via CoreML) runtimes.

\paragraph{Cascade and Adaptive Inference.}
Cascade systems and early\hyp exit networks~\cite{teerapittayanon2017branchynet,huang2018msdnet,han2024earlyexit_survey} offer an alternative efficiency strategy: route easy samples through lightweight classifiers and escalate hard ones to more expensive experts. Classical examples include Viola\hyp Jones face detection~\cite{viola2001rapid}, which uses a cascade of progressively complex classifiers to reject non\hyp faces early. In deep learning, BranchyNet~\cite{teerapittayanon2017branchynet} attaches shallow classifiers to intermediate layers, enabling early termination for confident predictions. Recent work on adaptive computation time (ACT)~\cite{graves2016act} and dynamic depth networks~\cite{wang2018skipnet} learns input\hyp dependent routing policies.

However, most cascade work focuses on maximizing accuracy or minimizing FLOPs on fixed test distributions, without exposing interpretable operational metrics (e.g., escalation rate, Stage~1 leakage rate) that platform operators need to tune cost--quality trade\hyp offs. Furthermore, cascade detectors for deepfake detection remain rare in the literature, and seldom report performance under distribution shift.

\paragraph{Calibration and Distribution Shift.}
Confidence calibration via temperature scaling~\cite{guo2017calibration,minderer2021revisiting} aligns predicted probabilities with empirical frequencies, improving reliability diagrams and expected calibration error (ECE). Well\hyp calibrated models are particularly valuable for cascade routing: thresholds set on calibrated probabilities better reflect true uncertainty and generalize more robustly to OOD data. However, calibration alone does not eliminate performance gaps under domain shift~\cite{pei2024deepfakesurvey}; a well\hyp calibrated model with 75\% accuracy on one distribution may still exhibit poor calibration if its true accuracy on another distribution is 60\%. Addressing cross\hyp dataset gaps fundamentally requires domain adaptation techniques~\cite{ganin2016dann,li2021dabn,liu2025unsupervised_da}, diverse training data, and transparent reporting of OOD performance.

\paragraph{Our Approach.}
Given these challenges, we propose a two\hyp stage cascade system for mobile deepfake detection that explicitly optimizes for both detection quality and deployment efficiency. Our Stage~1 filter uses MobileNetV4~\cite{mobilenetv4}—a lightweight CNN optimized for mobile hardware—to perform fast frame\hyp level screening. Samples with highly confident Stage~1 predictions (very likely real or very likely fake) receive immediate labels, while ambiguous cases escalate to a Stage~2 expert (EfficientNetV2\hyp B3~\cite{efficientnetv2} or optional GenConViT~\cite{genconvit_github}) for high\hyp fidelity classification. We formulate threshold selection as a cost\hyp sensitive optimization problem: minimize false negative rate (FNR) subject to a bounded Stage~2 escalation budget, exposing interpretable trade\hyp offs between recall and compute cost.

We train both stages on a combined multi\hyp dataset manifest (CelebDF\hyp v2, FaceForensics++, DFDC, DeeperForensics\hyp 1.0) to promote cross\hyp dataset generalization, experiment with hard example mining to focus Stage~2 capacity on cases where Stage~1 is uncertain, and validate the cascade on the in\hyp the\hyp wild Deepfake\hyp Eval\hyp 2024 benchmark~\cite{deepeval2024} to measure OOD performance. In our main results and public training script, we use the simpler standard configuration with uniform sampling and treat hard\hyp example mining as an ablation. For mobile deployment, we apply post\hyp training quantization (INT8) and export calibrated ONNX models with fixed thresholds, achieving sub\hyp 200\,ms latency on a Xiaomi~13 smartphone (Snapdragon 8 Gen~2).

The remainder of this paper proceeds as follows: Section~3 surveys related work in detail, positioning our contributions within the broader landscape of deepfake detection, cascade architectures, and mobile deployment. Section~4 describes our datasets and task formulation. Section~5 presents our six\hyp stage methodology, including formal problem definitions, cascade optimization, and hard example mining. Sections~6 and~7 detail experimental setup and results on in\hyp distribution validation and OOD benchmarks. Section~8 reports ablation studies, and Section~9 discusses mobile deployment specifics (ONNX export, on\hyp device latency, accuracy). Finally, Section~10 discusses limitations and future work, and Section~11 concludes.
